\section{Introduction}\label{sec:introduction}

A pulse threshold in a delay equation is an intersection problem in a space
of histories.  Starting from a quiet equilibrium, a one-unit forcing pulse of
amplitude \(J\) produces a release history and, after autonomous evolution, a
curve on a transverse section.  If an unstable periodic orbit has one
unstable section direction, its local stable set is a codimension-one sheet.
The threshold sought here is the value of \(J\) at which the event-aligned
pulse curve crosses that sheet.

This geometric definition separates three statements that are often conflated.
A transverse voltage event merely selects a common section.  A
pulse/stable-sheet intersection is a local invariant-manifold statement.  It
becomes an onset between two observable regimes only after trajectories on
the two sides are proved to enter different attraction tubes.  The present
paper proves neither the final intersection nor this two-sided routing; it
organizes and validates the return-map estimates that reduce the remaining
work to explicit inequalities.

Delayed FitzHugh--Nagumo oscillations and delayed canards have been studied by
\citet{campbell2009delay,krupa2016complex,krupa2016explosion}; broader
slow--fast geometry is surveyed in
\citet{wechselberger2012canards,desroches2012mixed}.  Standard RFDE theory
provides semiflows and local Poincar\'e maps
\citep{hale1993introduction,diekmann1995delay}, while collocation and
pseudospectral methods provide effective numerical access to periodic
solutions and their spectra
\citep{engelborghs2001collocation,engelborghs2002ddebiftool,
ando2020collocation,dewolff2021pseudospectral}.  The obstruction addressed
here is more specific: a graph-transform proof on a ball of arbitrary
continuous histories needs event-aligned first and second derivatives in the
complete-history norm.  Finite sampling alone cannot supply that operator
bound, and separately norming the unstable direction and its adjoint
deflation destroys the cancellation that makes the stable block small.

The proved mechanism has four parts.  First, a selected event whose window
begins after two maximum delays gives a \(C^2\) complete-history return.
Second, a correlated adjoint quotient and a signed-bimeasure residual retain
the cancellations required by the stable output.  Third, a direct
two-period calculation transfers the hyperbolic splitting and proves the
stable and unstable-inverse rates without assuming that the nonlinear map is
the square of a one-period return.  Fourth, a source-bound pulse calculation
gives one common event branch, its moving-time derivative, opposite endpoint
functional signs, and an explicit stable-coordinate ball.

These ingredients do not yet form a threshold theorem for the concrete model.
The selected return is currently proved only on a ball smaller than the
preferred graph domain; none of the six full-ball Hessian caps has been
validated.  Consequently the available stable graph is qualitative, its
identification with the periodic-orbit stable set has not been checked on the
needed domain, and the pulse crossing remains conditional.  On the outer
side, a local attracting return exists on a microscopic section ball, but the
post-pulse history has not been attached to it.  Stating these boundaries at
the outset is part of the mathematical result: a section event, a stable-sheet
crossing, and a routed onset are different objects.

\paragraph{Organization.}
\Cref{sec:model} defines the RFDE and the pulse-threshold geometry.
\Cref{sec:return-tools} gives the general selected-return and stable-set
lemmas and a conditional completion theorem.  The model-specific inner
return and graph estimates occupy \cref{sec:inner-return}; the physical pulse
event and conditional crossing are treated in \cref{sec:pulse-crossing}.
\Cref{sec:outer-return} records the outer local return and the unresolved
attachment, and \cref{sec:discussion} collects the exact remaining gates.

